\documentclass[a4paper, 12pt]{article}
\usepackage{fontspec}
\usepackage{xeCJK}
\usepackage[left=2cm, right=2cm, top=2cm, bottom=2cm]{geometry}
\usepackage{tcolorbox}
\usepackage{fancyhdr}
\usepackage{amsmath}
\pagestyle{fancy}
\lhead{高中物理}
\rhead{南昌学而思}
\cfoot{generated by Linghui Ding}
\newcommand{\defitem}[2]{
\noindent \textsf{#1} \hspace{1em} \itshape{#2} \\
}
\newcommand{\formula}[2]{
  \noindent \textbf{#2} \\ \[ #1 \]
}
\newtcolorbox{arcbox}{colframe=orange, arc=3mm, boxrule=1pt}
\newcommand{\noticebox}[1]{
  \begin{arcbox}
    \normalfont #1
  \end{arcbox}
}
\begin{document}
\defitem{机械运动}{物体的空间位置随时间的变化}
\defitem{质点}{有质量的点用于替代物体。物体的大小、形状和体积对所研究问题的影响可以忽略不计}
\defitem{参考系}{处理和使用参考系满足四个性质}
\defitem{标准性}{作为参考系的物体被假定是不动的，并且其他的物体的运动与静止都需要相对于这个物体}
\defitem{任意性}{参考系的选择是任意的。但是处理问题时一般不能选择物体本身}
\defitem{统一性}{研究多个物体需要选同一个参考系}
\defitem{差异性}{同一个物体的运动，在不同参考系下一般是不同的}
\defitem{时间}{时间间隔与一段进行的过程相关；时刻是特定的一瞬间}
\defitem{路程}{运动轨迹的长度}
\defitem{位移}{从初位置指向末位置的有向线段}
\defitem{速度（瞬时速度）}{$v=\lim_{\Delta t\to0}\frac{\Delta x}{\Delta t}$ 速度是真正的描述物体运动以及快慢的物理量，作为一个矢量，方向代表运动的方向，大小代表那一时刻的运动快慢}
\defitem{平均速度}{$\bar{v}=\frac{\Delta x}{t}$ 位移除以时间，虽然也是矢量，但其方向不具有任何实际物理意义，大小用来近似表示速度的大小，只能反应一段时间内的快慢水平}
\defitem{瞬时速率}{$v_\text{速率}=|v|$ 瞬时速度的大小，由于速度是矢量，所以瞬时速率是一个标量，并且大于等于0}
\defitem{平均速率}{$\bar{v_\text{速率}}=\frac{s}{t}$ 平均速率是路程与时间的比值，能反应日常生活中的运动快慢，与运动轨迹相关}
\defitem{x-t图像}{
  图像中直线斜率代表速度。曲线的切线斜率代表速度。\\
  斜率为负的意味着速度方向和拟定的正方向相反。\\
  图线与时间轴的交点代表物体在坐标为0的位置，图线与纵坐标轴的交点代表物体在初时刻的位置。\\
  图像中任意两个点的\underline{连线斜率}代表两个时刻间的平均速度。
}
\defitem{加速度}{
  \[a=\frac{\Delta v}{t}\]
  加速度也是一个矢量，方向代表速度变化的方向，大小代表速度变化的快慢。\\
  加速度与速度变化量的方向相同。下图是速度变化量的求解方法，需要讲速度的起点对齐，然后从初速度的末尾指向末速度的末尾。\\
  加速度和速度完全没关系。就好像有人问你在哪？你说我刚出门走了200米一样。\\
  和速度类似，不存在加速度与速度变化量成正比这样的说法，只有当a一定时，速度变化量与时间成正比。\\
  由于匀速直线运动，速度不随时间发生变化，所以加速度为0。 \\
  a与v同向，物体做加速运动。\\
  a与v反向，减速运动到0，后反向加速运动。\\
  a越大，速度变化越快，a越小，速度变化越慢。
}
\noticebox{
  你知道吗，上面的加速度定义其实是平均加速度的定义式，因为只有在平均物理量，我们才会使用一段时间，并且不要求时间很短。我们的瞬时加速度定义式有两种，一个是和瞬时速度类似的，另一个是利用牛顿第二运动定律（在后期课程大家才会接触）。如下，
  \[a=\lim_{t\to0}\frac{\Delta v}{t},a= \frac Fm\]
}
\defitem{a-t图像}{
  这个图像里面，图像与时间轴包围的面积代表这段时间内的速度变化量。\\
  图像在时间轴上方，那么面积认为是正的，否则就是负的，正的面积代表速度变化量为正。\\
  图像与时间轴的交点意味着此时的加速度为0，一般而言加速度为0，左右如果图像分布在时间轴的上下两侧，物体运动减速，注意仅仅是减速。
}
\defitem{匀变速直线运动}{为加速度恒定的直线运动，大小和方向都不变。所以有下面的公式成立：
  \[a=\frac{v-v_0}{t}=\frac{\Delta v}{t}\]
}
\defitem{v-t图像怎么用？}{匀变速直线运动的图像在v-t图像中，表达式为$v=at+v_0$，也就是斜率是a，截距是$v_0$的直线。}
\defitem{匀变速直线运动x-v公式}{
  \[v^2=v_0^2+2ax\]
}
\defitem{自由落体运动}{
  就是特殊的匀加速直线运动，初速度为0，加速度为重力加速度，方向竖直向下。下面是基本公式：
  \[x=\frac12gt^2\]
  \[v=\sqrt{2gx}\]
  \[t=\sqrt{\frac{2x}{g}}\]
}
\noticebox{
  你知道吗？自由落体类似的匀变速直线运动有着非常有趣的二级结论：
  \begin{itemize}
    \item 相同时间间隔每段位移之比是1:3:5
    \item 相同位移间隔每段时间之比是1:$\sqrt2-1$:$\sqrt3-\sqrt2$
  \end{itemize}
}
\defitem{竖直上抛运动}{
  就是特殊的匀减速直线运动，并且可能包含返回。\\
  初速度不等于0，加速度为重力加速度，方向竖直向下。公式只需要把匀变速直线运动的加速度替换即可，此处省略。\\
  比较有趣的事情是，匀变速直线运动中，所有的返回加速过程与减速过程是关于返回点对称的，也就是说物体减速经过某个点，之后再返回过程中将会是速度大小相同，但是方向相反的。
}
\noticebox{
  运动学的多解，可以是位移大小相同，方向未知，也可以是速度大小相同，方向未知。甚至可以是两者同时存在。使得我们会有两解、甚至三解的情况。
}
\defitem{5个基本公式}{
  \[v=v_0+at\]
  \[x=v_0t+\frac12at^2+x_0\]
  \[x=vt-\frac12at^2+x_0\]
  \[x=\bar vt=\frac{v_0+v_t}{2}t\]
  \[v_t^2-v_0^2=2ax\]
  在诸多公式里面我最喜欢，也最擅长使用两个。\[x=\frac{\bar v}{t}=\frac{v_0+v_t}{2t}\]这个公式也有一个文字描述，中间时刻的瞬时速度等于平均速度。只要发现题目给了一段位移并且有这段位移所花费的时间，那么这个公式是首选。然后计算出中间时刻的速度。\[v^2-v_0^2=2ax\]这个公式由于缺少时间，所以特征非常明显，一般情况下，如果题目缺少很多时间变量，那么可以考虑用这个公式。
}
\defitem{x-t图像理解}{
  \emph{（切线）斜率代表速度}，正斜率表示正向运动，负斜率表示反向运动。

  \emph{在图像的转折点或者是v=0，并且左右斜率正负不相同的地方}是运动反向点。

  匀速直线运动是任意一条水平线。

  匀变速直线运动是一条一元二次函数图像。
}
\defitem{v-t图像理解}{
  \emph{（切线）斜率代表加速度}，正斜率表示加速运动，负斜率表示减速运动。

  \emph{与t轴的交点，左右如果正负不同}，那么这就是运动反向点。

  匀速直线运动是一条$v=0$的水平线。

  匀变速直线运动是一条倾斜的直线。

  与t轴围成的面积代表位移，面积在t轴上方表示正向位移，面积在t轴下方表示反向位移。

  \underline{v-t图像没有位置信息，只有位移，所以无法直接确定多物体的相对位置关系。}
}
\defitem{拓展公式}{
  \[v_{\frac12t}=\bar v=\frac{v_0+v_t}{2}\]
  \[v_{\frac12x}=\sqrt{\frac{v_0^2+v_t^2}{2}}\]
  一般情况下，相等时间间隔$T$，第M段与第N段的位移之差为
  \[\Delta x_{M,N}=(M-N)aT^2\]
  在初速度为0时有两个特殊结论：下面是相等时间间隔，
  \[\text{累计速度比}=1:2:\cdots:n\]
  \[\text{累计位移比}=1:2^2:\cdots:n^2\]
  \[\text{每段位移比}=1:3:\cdots:2n-1\]
  下面是相等位移间隔，
  \[\text{每段时间比}=1:\sqrt{2}-1:\cdots:\sqrt{n}-\sqrt{n-1}\]
  \[\text{累计速度比}=1:\sqrt{2}:\cdots:\sqrt{n}\]
  \[\text{累计时间比}=1:\sqrt{2}:\cdots:\sqrt{n}\]
}
\formula{G=mg}{重力}
\defitem{重力}{
  方向常考，竖直向下，不是垂直向下，可以是垂直水平面向下。\\
  重力不会因物体的运动而改变，只有两个因素，质量和重力加速度。重力加速度本身也只和物体所在位置和海拔高度。\\
  重力的产生原因：万有引力+圆周运动。\\
  重力是最简单的力，所以分析的时候一般先分析重力。
}
\defitem{弹力}{有形变，想要恢复形变的力叫弹力，弹力的方向：
  一般情况下，有明确的接触面的弹力，方向沿着接触面法线方向。（就是垂直于接触面指向受力物体）\\
  但是接触面很多种：面面接触时接触面法线方向；点面接触时面的法线方向；点点接触时分为有无圆面（因为圆一般和其他物体接触只有一个接触点），有圆面一定是圆的切面法线方向指向受力物体，否则就是两个点的重心的连线方向。\\
  绳子提供的弹力只能沿绳方向，只能提供拉力。杆子可以像绳子一样提供拉力或支持力，但是杆子的弹力可以是沿杆方向也可以是任意方向，取决于杆子的一端是否固定。\\
  发生弹性形变时，弹簧的弹力与弹簧的形变量成正比。比值叫做劲度系数。劲度系数只与弹簧本身有关，和弹簧的圈数呈正比，和材料有关，与直径有关，等等。\\
  \[x=|l-l_0|\]
  上面的形变量的计算公式，注意里面的$l_0$指弹簧不受力时长度，叫做原长。绝对值即弹簧处于压缩状态，弹簧长度将会小于原长。弹簧处于压缩状态，弹力向两侧；弹簧处于拉伸状态，弹力向内。弹簧弹力等于任何一端的受力。弹簧两端的受力一定大小相等，方向相反，不管处于何种运动状态下。弹力只和形变量有关，也就意味着，弹力不会直接因为物体运动状态改变而改变，除非弹簧形变量变化。
}
\defitem{弹簧的劲度系数}{
  并联弹簧劲度系数 $k_\text{并} = k_1 + k_2 + \cdots + k_n$
  串联弹簧劲度系数 $\frac1{k_\text{串}} = \frac{1}{k_1} + \frac{1}{k_2} + \cdots + \frac{1}{k_n}$
  当只有两根弹簧时，串联和并联的关系可以简化为：
  \[
    k_\text{并} = k_1 + k_2, \quad k_\text{串} = \frac{k_1 k_2}{k_1 + k_2}
  \]
  由此可以知道，\underline{把一根弹簧进行分割时，如果分成两段，分割后两段的劲度系数是原来的两倍！}
  分成$n$段，分割后每段的劲度系数是原来的$n$倍。
  之后并联这些分段，整体的劲度系数达到了惊人的$n^2$倍！
}
\formula{f_{\text{滑}} = \mu N}{滑动摩擦力}
\formula{f_{\text{静}} \leq  f_{\text{最大静摩擦力}}}{静摩擦力}

\defitem{摩擦力定义}{摩擦力是两个\emph{相互接触}并有\emph{相对运动趋势或相对运动}的物体之间，在接触面上产生的阻碍相对运动的力。}

\defitem{摩擦力方向}{摩擦力的方向总是与相对运动方向相反，或与相对运动趋势方向相反。}

\defitem{摩擦力分类}{
  \emph{静摩擦力}：当物体间没有发生相对滑动，但有相对滑动趋势时，接触面之间产生的阻碍相对运动趋势的力叫静摩擦力。静摩擦力的大小随外力变化而变化，其最大值称为最大静摩擦力。
  \emph{滑动摩擦力}：当物体间发生相对滑动时，接触面之间产生的阻碍相对滑动的力叫滑动摩擦力。滑动摩擦力的大小通常为恒定值。
  \emph{滚动摩擦力}：当物体在另一物体表面滚动时，产生的阻碍滚动的力叫滚动摩擦力。滚动摩擦力远小于滑动摩擦力。
}

\defitem{摩擦力产生条件}{必须有两个物体相互接触，
  相互产生弹力，
  并且接触面具有粗糙性（微观不平），
且有相对运动或相对运动趋势。}

\defitem{摩擦力方向详述}{总是与物体相对运动方向或相对运动趋势方向相反。}

\defitem{摩擦力计算公式}{最大静摩擦力与正压力成正比，而静摩擦力本身与正压力无关。滑动摩擦力：$f = \mu N$，其中$\mu$为滑动摩擦因数，$N$为正压力。}

\defitem{摩擦因数}{摩擦因数与接触面的材料和粗糙程度有关，与接触面积大小无关。}

\defitem{静摩擦力特点}{静摩擦力的大小不一定等于最大静摩擦力，而是根据实际需要"自适应"变化，直到达到最大静摩擦力为止。一旦外力超过最大静摩擦力，物体就会发生滑动，转为滑动摩擦力。}

\noticebox{
  判断相对运动或相对运动趋势的方法有以下几种：
  \defitem{速度法}{判断两个物体在接触面方向上的速度是否相同。如果速度不同，则存在相对运动；如果速度相同，则无相对运动，但需进一步判断是否有相对运动趋势。}
  \defitem{受力分析法}{分析外力作用下，若没有摩擦力或摩擦力不足，物体是否会发生相对滑动。若会，则存在相对运动趋势。}
  \defitem{极端假设法}{假设接触面光滑（无摩擦），判断物体是否会发生相对滑动。如果会，则存在相对运动趋势。}
  \defitem{几何关系法}{通过分析物体的运动轨迹、转动或连接方式，判断接触面是否会产生相对滑动或滑动趋势。}
  常用的判断顺序是：先判断有无相对运动（速度是否相同），若无则判断有无相对运动趋势（极端假设法或受力分析法）。
}

\defitem{力的合成与分解}{
  力的合成是指将两个或多个力合成一个等效的力（合力），使其对物体的作用效果与原来各力共同作用的效果相同。

  力的分解是指将一个力分解为两个或多个分力，使这些分力的合力等于原来的力。
}

\defitem{两力合成的大小范围}{
  已知两个力的大小分别为$F_1$和$F_2$，夹角为$\theta$，使用余弦定理得到合力$F$的大小为
  \[
    F = \sqrt{F_1^2 + F_2^2 + 2F_1F_2\cos\theta}
  \]
  合力的大小范围为
  \[
    |F_1 - F_2| \leq F \leq F_1 + F_2
  \]
  其中，当$\theta=0^\circ$时，合力最大$F=F_1+F_2$；当$\theta=180^\circ$时，合力最小$F=|F_1-F_2|$。
}

\defitem{超过两个力的合成}{
  对任何两个力先进行合成，得到一个合力，然后再将这个合力与下一个力进行合成，依次类推，直到所有力都被合成。
}

\defitem{分解为超过两个力}{
  对一个力进行分解时，可以选择任意两个方向进行分解，得到两个分力。然后可以将其中一个分力继续分解为两个方向的分力，依次类推，直到分解为所需的分力数量。
}

\defitem{多个力合成时大小范围}{
  任意两个力先求出范围$[|F_1 - F_2|, F_1 + F_2]$。\\
  接着第三个力与这个范围进行合成，需要注意最小值。举几个例子：\\
  三个力大小分别是4、4、6，那么首先得到$[0,8]$，接着得到[0,14]，其实就是力的大小范围不能是负数，而且6在范围内，所以可以得到最小的0.\\
  如果是4、4、10，那么首先得到$[0,8]$，接着，你会注意到我们无法获得0这个最小值，范围是$[2,18]$。\\
  如果是4、8、3，首先得到$4,12$，接着无法得到0，只能找最小值，于是有范围$1,15$\\
  \underline{上面的范围计算，虽然我们选取的顺序是特定的，但是结论是通用的，你不管怎么调整顺序不会影响结论。}
}

\defitem{合力的方向}{
  \underline{合力的方向总是在两个分力的夹角之间}。只有当两个力共线且方向相反时，合力才可能与其中一个分力方向相同或相反。一般情况下，合力一定落在两个分力的夹角内。
}

\formula{\sin^2\theta + \cos^2\theta = 1}{三角恒等式}

\defitem{余弦定理与正弦定理}{
  \textbf{余弦定理}：\\
  任意三角形，设三边为$a,b,c$，对应的对角为$A,B,C$，则有
  \[ c^2 = a^2 + b^2 - 2ab\cos C \]
  \textbf{正弦定理}：\\
  \[ \frac{a}{\sin A} = \frac{b}{\sin B} = \frac{c}{\sin C} \]
}
\defitem{将来大家会用到的公式}{
  \[ a\sin\theta + b\cos\theta = \sqrt{a^2 + b^2}\sin(\theta + \alpha) \]
  其中，$\tan\alpha = \frac{b}{a}$。
}
\defitem{受力分析怎么做？}{重力先画，不一定要画在重心位置，或者是中心点。\\
  挑一个好看的地点更好，比如图中已经有很多直线了，选择那些直线的交点，这样可以让图像更清晰。\\
  弹力是一个难点，牢记我们之前的规律应用他们，\\
  最后再画摩擦力，因为有弹力才能有摩擦力。\\
  总结为1重2弹3摩擦。最后题目中还有别的外力再补充。

  二力平衡：一定是等大反向，作用于同一个物体，方向在一条直线上。

  三力平衡：不共线的三力平衡一定是形成一个首尾相连的密闭三角形。

  多个力平衡：任意一个力与其他力的合力是二力平衡。注意多个力的合力不要和其他任何一个分力同时存在，一旦选择了合成，那么这些力就应该消失。

  正交分解处理平衡：只需要在两个垂直的方向上，合力为0就是平衡。三个力及以下，我们很少使用正交分解方法，最多使用正弦定理。四个力则基本上都要用正交分解解决。
}

\defitem{正交分解求解共点力平衡}{
  共点力平衡问题常用正交分解法。将所有力分解到两个互相垂直的方向（如水平和竖直，或沿斜面和垂直斜面），分别列出每个方向上的分力代数和为零的方程：
  \[
    \sum F_x = 0,\quad \sum F_y = 0
  \]
  这样可以将力的平衡问题转化为两个独立的代数方程，便于求解未知力的大小和方向。
}
\defitem{绳子模型}{
  绳子只能传递拉力，方向总是沿绳子本身。
  \emph{活结}：绳子可以自由滑动，绳子的拉力在不同地方相等。
  \emph{死结}：绳子固定，不同段的力的大小可以不相同。
}

\defitem{杆子模型}{
  杆子可传递拉力和压力。
  \emph{自由杆}：两端都不固定，杆内力沿杆方向，通常只考虑拉力或压力。一般使用"铰链"链接。
  \emph{固定杆}：一端或两端固定，提供的力可以是任意方向，任意大小，一般而言一定会强调"固定在墙壁上"。
}
\defitem{三个力的动态平衡问题}{
  情形一、重力恒定，另一个力方向恒定，第三力变化。设$G$为重力，$F_1$方向恒定，$F_2$大小可以发现变化过程为先变小后变大，如果$F_1$的与重力的夹角是$90^\circ$那么，只有变小过程。

  情形二、重力恒定，另一个力大小恒定，第三力变化。设$G$为重力，$F_1$大小恒定，则利用余弦定理知道，当两个力的大小不变的时候，随着$F_1$与重力的夹角从小变大的过程，$F_2$在逐渐增大。
  设$G$为重力，$F_1$大小恒定，则利用余弦定理知道，当两个力的大小不变的时候，随着$F_1$与重力的夹角从小变大的过程，$F_2$在逐渐增大。
}
\defitem{伽利略理想斜面实验}{
  伽利略通过让小球在不同倾角的光滑斜面上滚动，发现小球在下斜面上速度会逐渐增加，而在上斜面上速度会逐渐减小。\\
  如果让小球从一个斜面上滚下，再沿另一个斜面上升（两斜面高度相同且光滑），小球总能回到原来的高度。\\
  伽利略进一步推理：如果上升的斜面无限延长且完全光滑，小球将永远以恒定速度运动下去，不会停止。

  这一实验说明：力不是维持物体运动的原因，而是改变物体运动状态（速度大小或方向）的原因。物体之所以最终停止，是因为受到了摩擦力等阻力的作用。如果没有外力，物体将保持匀速直线运动。这为牛顿第一运动定律（惯性定律）的提出奠定了实验基础。
}

\defitem{牛顿第一运动定律（惯性定律）}{
  质量是惯性的量度，质量越大惯性就越大。惯性是物体改变运动状态难易程度的体现。为什么我们之前质点，保留了质量呢？就是因为质量会影响惯性，而惯性会影响物体运动状态改变。所以运动学问题需要质量。

  惯性参考系、将相对于地面这样的参考系，静止或者做匀速直线运动的参考系称为惯性参考系（简称惯性系）。在惯性系中牛二是成立的。

  非惯性参考系、所有相对于地面这样的参考系，有着不等于0的加速度的参考系都是非惯性参考系（简称非惯性系）。在非惯性系中，牛顿运动定律不再适用，需要引入虚拟力来分析物体的运动。这个虚拟力叫做非惯性力，大小等于物体的质量乘以参考系的加速度，方向与参考系的加速度方向相反。
}

\defitem{牛顿第三运动定律（相互作用力）}{
  \[F_{12}=-F_{21}\]
  作用力与反作用力大小相等，作用物体和施加物体不同，方向相反，作用于不同物体。
}

\formula{F_\text{合}=ma}{牛顿第二运动定律}

\defitem{牛顿第二运动定律（重要定律）}{
  成功联系运动和力，并且知道运动是力的反映，力可以影响运动。

  公式中的加速度是瞬时加速度，这才是加速度的定义式，能够决定任意时刻的加速度。并且只要物体的力改变了，物体的加速度就会改变。

  静力平衡是加速度为零的情况。对于之前我们的正交分解多了一种分解方法，沿加速度和垂直加速度分解，这样就不用分解加速度。垂直方向的力之间互不影响，有着各自的牛顿运动定律表达式。
}

\defitem{国际单位制（量纲法）}{
  这是一个高考考点，简单来说就是单位可以计算。做选择题至此大家已经收获了一个强大的方法：量纲法，选择题做之前先用量纲法试一试。

  记忆方法：7个单位，m、kg、s、A、K、mol、cd。前三个是力学的，接着是电磁学的电流单位，热学两个（温度和物质的量），最后一个发光强度。
}

\noticebox{
  高中阶段一般默认地面做参考系。图中曲线轨迹为运动轨迹，曲线的长度就是轨迹的长度也是路程。位移变化量和位移的区别是位移变化量的起点和终点随着每次运动而变化，位移的起点只有一个。后续学习中，位移都是相对于一个固定的初始位置的，末位置是任意时刻的位置。路程是长度，是一个值，不是轨迹。

  你会注意到，上面是x关于t的一元二次方程，注意加速度里面有$\frac12$。上面两个公式，是通过两个不同的角度观察一元二次方程，从初始分析，和从结束逆推。
}

\end{document}
